\section*{Supplementary Information}

\subsection*{S1: Mesh Convergence and Diffusion-Layer Resolution}
\label{sec:s1}

The optimisation was repeated on four structured meshes:
$20\times5$, $40\times10$, $80\times20$ (primary), and $160\times40$
(validation), all for the linear gradient target with 400 iterations.

\begin{table}[h]
\centering
\caption{Mesh convergence for the linear gradient target.
         The diffusion layer $\delta\approx4\,\mu$m is underresolved
         at all resolutions; RMSE measures the global outlet profile
         error, not the local diffusion-layer accuracy.}
\label{tab:mesh_conv_full}
\begin{tabular}{@{}lcccl@{}}
\toprule
Mesh & Elements & $\Delta y\,(\mu\text{m})$ & RMSE & Note \\
\midrule
$20\times5$   & 100  & 100  & 0.584        & Screening only \\
$40\times10$  & 400  & 50   & 0.211        & \\
$80\times20$  & 1600 & 25   & \rmseTopOpt  & Primary result \\
$160\times40$ & 6400 & 12.5 & \rmseHiRes   & Validation \\
\bottomrule
\end{tabular}
\end{table}

The topology (branching structure, number of junctions) is qualitatively
identical between the $80\times20$ and $160\times40$ meshes, confirming
that the channel-scale structure is mesh-converged at the primary
resolution.
The remaining RMSE at $160\times40$ reflects the unresolved diffusion
sublayer ($\delta/\Delta y \approx 0.32$ at the finest level); adaptive
mesh refinement targeting the fluid--solid interface is deferred to
future work.

For larger meshes, the direct MUMPS solver is replaced by a
block-triangular field-split preconditioner: algebraic multigrid (HYPRE)
for the velocity block and the Schur complement approximation for the
pressure~\cite{elman2014}.
Iteration counts remain below 50 for all mesh sizes tested
($20\times5$ to $160\times40$), confirming near-optimal scalability.

\subsection*{S2: Filter Radius Sensitivity}
\label{sec:s2}

The Helmholtz filter radius $r_f = 2\Delta x$ was varied between
$1\Delta x$ and $4\Delta x$ on the $80\times20$ mesh.
The outlet RMSE varied by less than $0.02$ across this range, and
the qualitative topology was preserved throughout.
Filter radii below $1\Delta x$ produced checkerboard artefacts;
radii above $4\Delta x$ over-smoothed the design and increased RMSE.
The value $r_f = 2\Delta x$ lies in the centre of the stable range
and is used throughout the main paper.

\subsection*{S3: SUPG Stabilisation Validation}
\label{sec:s3}

Three forward solves on the $80\times20$ mesh were compared for the
linear gradient target:
(i)~no stabilisation, (ii)~SUPG with $\tau$ from Eq.~\eqref{eq:supg},
and (iii)~full upwinding ($\tau = h/(2\|\mathbf{u}\|)$).
Without stabilisation, the concentration field exhibits oscillations with
amplitude up to~$0.4$ at the inlet interface.
SUPG reduces oscillation amplitude below $10^{-3}$ while maintaining
the smooth outlet gradient profile.
Full upwinding eliminates oscillations but introduces excessive numerical
diffusion, reducing the outlet gradient slope by approximately~$30\%$.
SUPG is therefore the appropriate choice for this operating regime.

\subsection*{S4: OC vs.\ MMA Convergence Comparison}
\label{sec:s4}

Both the Optimality Criteria (OC) method and the Method of Moving
Asymptotes (MMA)~\cite{svanberg1987} were run for 400 iterations on the
linear gradient target ($80\times20$ mesh, identical initialisation and
continuation schedules).
OC converged to a final objective of $J = \finalJ$ with outlet
RMSE~$= \rmseTopOpt$.
MMA converged to a comparable objective (within~$5\%$ of the OC result)
but required approximately $1.4\times$ the wall-clock time due to
additional inner-loop solves.
The optimised topologies from the two methods are qualitatively similar
(same branching structure), confirming that the result is not specific
to the OC update rule.
OC was used for all main-paper results.

\subsection*{S10: Adjoint Gradient Taylor-Test Details}
\label{sec:s10}

The Taylor remainder test (Section~\ref{sec:taylor_test}) was performed
on the $20\times5$ mesh for computational speed.
A random unit-norm perturbation $\delta\rho$ was drawn from a
standard-normal distribution (seed~$= 42$) and normalised in the $L^2$
norm.
The first-order Taylor remainder $R(\varepsilon)$ was evaluated for
$\varepsilon \in \{10^{-8}, 10^{-7}, \ldots, 10^{-2}\}$.

The least-squares slope over $\varepsilon \in [10^{-8}, 10^{-4}]$
is $2.01 \pm 0.03$, confirming:
\begin{itemize}
    \item the adjoint equations~\eqref{eq:adj_flow_mom}--\eqref{eq:adj_conc_bc}
          are correctly implemented;
    \item the chain rule through the Helmholtz filter (Eq.~\eqref{eq:helmholtz})
          and Heaviside projection (Eq.~\eqref{eq:heaviside}) is applied
          exactly;
    \item the SUPG stabilisation of the adjoint transport equation does not
          corrupt the gradient at this level of precision.
\end{itemize}
The test is executed automatically on every commit via the CI pipeline
(\texttt{python -m micrograd.taylor\_test}).
\subsection*{S5: Thresholding Robustness and $D_{\min}$ Sensitivity}
\label{sec:s5}

To assess porous leakage, the converged near-binary design ($\beta=64$,
grey-zone fraction $<2\%$) is hard-thresholded at $\bar\rho=0.5$ and
the forward Brinkman problem is re-solved with $D_{\min}$ retained in
solid regions.
The analytical bound $t_{\mathrm{diff}}/t_{\mathrm{flow}} \approx 430$
(with $p=3$, $D_{\min}=10^{-15}\,\text{m}^2\,\text{s}^{-1}$) confirms
that diffusive transport through solid is negligible in steady state.

\begin{table}[h]
\centering
\caption{Binary validation results: RMSE and hydraulic resistance
         for the grey Brinkman, binary Brinkman, and binary Stokes designs
         on the same $80\times20$ mesh.}
\label{tab:binary_val}
\begin{tabular}{@{}lccc@{}}
\toprule
\textbf{Design} & \textbf{RMSE} & \textbf{$\Delta$RMSE vs grey} & \textbf{Method} \\
\midrule
Grey Brinkman  & \rmseTopOpt         & ---         & Smooth $\alpha(\bar\rho)$ \\
Binary Brinkman & \rmseBinaryBrinkman & $<1\%$      & $\bar\rho\in\{0,1\}$, smooth $\alpha$ \\
Binary Stokes   & \rmseBinaryStokes   & $<1\%$      & $\bar\rho\in\{0,1\}$, $\alpha=0$ in fluid \\
\bottomrule
\end{tabular}
\end{table}

\noindent
The sub-$1\%$ RMSE change between the grey Brinkman design and the
binary Stokes design (where $\alpha=0$ exactly in all fluid elements,
eliminating the Brinkman penalisation entirely) confirms two things:
\begin{enumerate}
    \item The optimiser does \emph{not} exploit grey-zone diffusion
          pathways; the near-binary design ($\beta=64$, grey fraction
          $<2\%$) performs identically to the hard-thresholded binary.
    \item The Brinkman penalisation does not introduce a significant
          velocity artefact in fluid regions; setting $\alpha=0$
          (free Stokes flow) gives the same outlet concentration profile.
\end{enumerate}
These results directly address the concern that the Brinkman surrogate
may produce non-physical concentration gradients that disappear upon
thresholding.
The $t_{\mathrm{diff}}/t_{\mathrm{flow}} \approx 430$ analysis
(Section~\ref{para:leakage}) is consistent with this numerical evidence:
diffusive leakage through solid regions is negligible in steady state.

